\begin{problem}{Shugou of the Year}{standard input}{standard output}{10 seconds}{256 megabytes}

佐賀大学は、今年の大学祭の期間に、学生たちが持ち寄って飲んだお酒と、その学生たちの所属学部情報を入手しました。この情報を適切に整理し、摂取したアルコールの総量が最も多い学部の名前を特定してください。与えられたクラスに宣言されたメソッドを適切に実装し、STLを活用する必要があります。


\InputFile
1行目に、お酒を持ち寄って飲んだ学生の数 N が与えられます。$(1 \le N \le 10^6)$
続くN行にわたり、各学生の所属学部 S、お酒の容量 V、アルコール度数 P がスペース区切りで与えられます。$(1 \le |S| \le 25, 1 \le V \le 1000, 0 \le P \le 100)$

Sはアルファベットの小文字と \tt{_} (アンダースコア)のみで構成される文字列で、常に以下のいずれかであることが保証されます。

- \tt{education}

- \tt{science_and_engineering}

- \tt{economics}

- \tt{agriculture}

- \tt{art_and_regional_design}

$V$と$P$は整数です。


\OutputFile
摂取したアルコール総量が最も多い学部の名前を出力してください。総量が最多の学部が複数ある場合は、辞書順で最も先に来る学部の名前を出力してください。



\end{problem}

